
\label{sec:prelim}

VLAs are trained by starting from a pre-trained vision-language model (VLM) backbone, and adapting it for robotic control. The training dataset $\data$ contains robot trajectories, which are sequences of observations $\bo_t$ and actions $\ba_t$. The observations $\bo_t = [\bI_t^1, \ldots, \bI_t^n, \bq_t]$ consist of $n$ camera images $\bI_t^i$ and the joint configuration of the robot $\bq_t$, while the actions $\ba_t$ consist of joint or end-effector commands. 

VLAs are typically trained to predict an \emph{action chunk}, corresponding to a short trajectory of future actions $\ba_{t:t+H}$, based on a recent history of observations $\bo_{t-T:t}$ (often a shorter horizon of actions, $\hat{H} < H$, is executed).
The action chunk can be generated by an ``action expert'', a smaller transformer that attends to the VLM backbone and enables fast inference at runtime. The action expert typically uses a flow matching~\citep{lipman2022flow} (or diffusion) objective that captures the multi-modality of the robot actions. To learn effective representations, our model also uses the knowledge insulation (KI) training recipe~\cite{driess2025knowledge}: the VLM backbone is supervised with FAST tokens \cite{pertsch2025fast}, and while the action expert attends to all of the activations in the VLM backbone, gradients from the action expert do not flow into the VLM backbone, such that the VLM is trained via the comparatively stable discrete cross-entropy loss.

In addition to the observation and action, each training example for the VLA is accompanied by 
a \emph{prompt} or 
\emph{context}, which we denote with $\mathcal{C}_t$. Conventionally this corresponds to a language instruction $\lang_t$, such that $\mathcal{C}_t = (\lang_t)$, provided by a human annotator (e.g., ``clean up the kitchen'').

In designing \MethodName{}, we explore how additional information added to the context for each training example can enable learning from diverse and heterogeneous datasets (including suboptimal behaviors and failures). As we will show, training with this data leads to a model with greater robustness and dexterity, and makes it possible for the model to generalize more broadly. The training objective for the VLA $\pi_\theta$ corresponds to an approximate log-likelihood given by
\begin{equation}
\max_{\theta} \; \E_{\data}
\left[\log \pi_{\theta}(\ba_{t:t+H} \mid \bo_{t-T:t}, \mathcal{C}_t)\right].
\label{eq:vla_il}
\end{equation}
Note that a flow matching action expert optimizes an approximate lower bound rather than a closed form log-likelihood \cite{black2024pi_0}. The dataset $\data$ typically consists of high-quality human demonstration trajectories; however, as mentioned, we use a broader dataset, which includes failed episodes and suboptimal autonomous rollouts, as well as other data sources such as egocentric human video data. We will show how using a sufficiently detailed and informative context $\mathcal{C}_t$ makes it possible to incorporate such diverse data and, perhaps surprisingly, even results in better policy performance and generalization.


\section{\texorpdfstring{\MethodName{}}{pi0.7} Overview}
\setcounter{figure}{1}
\begin{figure*}[t]
    \centering
    \includegraphics[width=\textwidth]{compressed_images_final/model_architecture.pdf}
    \caption{\textbf{Architecture overview.} The \MethodName{} model is a 5B-parameter VLA consisting of a 4B VLM backbone, a \texttt{MEM}-style video history encoder and a 860M parameter action expert. The model's context includes multiple distinct modalities, including language commands, episode metadata that describes the data quality and strategy, and multimodal inputs such as subgoal images. At runtime, the language commands are produced by a high-level semantic policy based on the same architecture, and the subgoal images are produced by a lightweight world model based on the BAGEL image generation model~\cite{deng2025emerging}.}
    \label{fig:arch}
\vspace{-10pt}
\end{figure*}

\MethodName{} is our newest robotic foundation model that builds on the existing VLA architecture from \Pizs{}~\cite{pi06model} and the \texttt{MEM} memory system~\cite{torne2026mem} and extends it with multi-modal context conditioning. The model consists of a VLM backbone initialized from the Gemma3 4B-parameter VLM \cite{gemmateam2025gemma3technicalreport} (including a 400M-parameter vision encoder),
and a flow matching action expert with 860M parameters. The model has about 5B total parameters. The vision encoder is also initialized from Gemma3 and follows the design of the \texttt{MEM} video history encoder~\cite{torne2026mem}, applying both temporal and spatial compression over history observations and outputting a fixed number of tokens for any number of history frames. An overview of the model architecture is provided in Fig.~\ref{fig:arch}, and Sec.~\ref{subsec:architecture} describes the architecture in more detail. 

Our previous models, \Piz{}, \Pizf{}, and \Pizs{}, use a short textual description of the task as the context. In training \MethodName{}, we expand the context to include additional information and modalities: more expressive language commands, episode metadata, and subgoal images, making it possible to train on diverse and potentially suboptimal data.


